\scanpage{15}

% ll. 1–2
\mspage{\textcircled{12}}
\uline{Proposition} $G$ groupe abélien (pour simplifier), $V=G\otimes_{\mathbb Z}\mathbb F_2$, $\varepsilon$ un $\mathbb F_2$ vectoriel, $\unsure{\Gamma|\mathscr T}=\Psi(G)=\operatorname{Ker}(G\times G\times G\xrightarrow[(x,y,z)\mapsto x+y+z]{}G)$,

% upper right, beside lines 2–5
\margin{$p_i:\Psi(G)\to G$ induit par $\mathrm{pr}_i$ $(1\leq i\leq3)$\\
$N_i=\operatorname{Ker}p_i\xleftarrow[\alpha_i]{\sim}G$\\
\deleted{$\unsure{g\mapsto(g,\widehat 0,g)|\{(x,\widehat x,x\}}$}}

% ll. 3–4
a) Pour tte extension centrale \deleted{\unsure{de|$E_0$}} :

% l. 4 displayed sequence
\[
1\longrightarrow\varepsilon\mathbin{-} E_0\longrightarrow G\mathbin{-}1
\]

% ll. 5–8
considérons $\Psi(E_0)=p_1^*(E_0)\wedge p_2^*(E_0)\wedge p_3^*(E_0)=\widetilde\Gamma$ (produit de Baer d'extensions de $\deleted{\unsure{\Psi_2}}\Gamma$ par $\varepsilon$, images inv de l'ext. $E_0$ (\deleted{\unsure{de}} $G$ par $\varepsilon$) par les $p_i$). Comme $p_j\alpha_i=\left\{\begin{array}{ll}0&\text{si }i=j\\ \mathrm{id}_G&\text{si }i\ne j\end{array}\right.$, on trouve que $\alpha_i^*(\widetilde E)$ est splitté canoniquement.

% ll. 9–10
d'où

% ll. 9–10, diagram
\diagram{\begin{tikzcd}[ampersand replacement=\&,column sep=1.6em,row sep=1.25em]
1 \arrow[r,dash]\& \varepsilon \arrow[r]\& \widetilde E \arrow[r]\& \Gamma \arrow[r,dash]\&1\\
\&\&\& G\arrow[ul,"\rho_i"]\arrow[u,"\alpha_i"']\&
\end{tikzcd}}

% ll. 11–12
Ceci posé, si
\[
q:\dot V\longrightarrow\varepsilon
\]

% ll. 13–14
est l'application quadratique qui définit la classe de l'extension $E_0$, on a pour $g\in G$

% l. 15
\[
\rho_1\deleted{\unsure{\alpha_1}}(g)\rho_{\unsure{2|1}}\deleted{\unsure{\alpha_2}}(g)\rho_3\deleted{\alpha_3}(g)=q(\dot g)
\]

% l. 16
où $\dot g\in V$ est la classe de $g\in\unsure{K|G}$ ds $V=G/2\unsure{C|G}$.

% ll. 17–23
b) Toute \inserted{syst. $(\widetilde E,\rho_1,\rho_2,\rho_3)$ d'une} extension \inserted{centrale} $E$ de $\Gamma$ par $\varepsilon$ et de relèvements $\rho_1,\rho_2,\rho_3$ de $\alpha_1,\alpha_2,\alpha_3$ provient d'une extension $E_0$ de $G$ par $\varepsilon$, déterminée à isom (non unique) près --- sa classe $q\in\mathrm{Quad}_{\mathbb Z}(G,\varepsilon)\simeq\mathrm{Quad}_{\mathbb F_2}(V,\varepsilon)$ est donnée par la formule précédente. La catégorie \inserted{(pour les iso)} des \deleted{triples} syst. $(\widetilde\Gamma,\rho_1,\rho_2,\rho_3)$ est \uline{rigide} i.e. tout autom est l'identité (donc tt autom. de l'ext. $E_0$ induit l'identité sur \inserted{l'ext} $\Psi(E_0)$ \dots.)
% archive sheet number: 14
